\documentclass{article}
\usepackage[UTF8]{ctex}
\usepackage{amsmath,amssymb,amsfonts}
\usepackage{enumitem}
\usepackage{geometry}
\geometry{a4paper, margin=1in}
% === 错题本排版优化（不参与题目标识） ===
\setlength{\parindent}{0pt}
\setlength{\parskip}{0.35em plus 0.10em minus 0.08em}
\setlength{\emergencystretch}{3em}
\setlist{itemsep=0.25em, topsep=0.40em, parsep=0pt}
\allowdisplaybreaks[2]
% === 错题本排版优化结束 ===
\begin{document}
\section*{中南大学2025年高等代数真题}
\section*{136．中南大学2025年高等代数试卷}

\subsection{题目1}
一、（16 分）设 \(p(x)\) 是 \(n\) 次有理系数多项式，\(n\) 为大于1的奇数，且 \(p(x)\) 在 \(\mathbb{Q}\) 上不可约。证明：如果 \(c_{1}, c_{2}\) 是 \(p(x)\) 的两个不同复根，那么 \(c_{1}+c_{2}\) 不是有理数．

\subsection{题目2}
二．（18分）求行列式

\[
\left|\begin{array}{cccc}
\sin \theta_{1} & \sin 2 \theta_{1} & \cdots & \sin 2025 \theta_{1} \\
\sin \theta_{2} & \sin 2 \theta_{2} & \cdots & \sin 2025 \theta_{2} \\
\vdots & \vdots & & \vdots \\
\sin \theta_{2025} & \sin 2 \theta_{2025} & \cdots & \sin 2025 \theta_{2025}
\end{array}\right|
\]

的值。

\subsection{题目3}
三．（18 分）设 \(f(x) \in C^{1}\left[x_{0}, x_{2}\right], x_{0}<x_{1}<x_{2}\) 。问是否存在次数不超过 2 的多项式 \(p(x)\) ，满足

\[
p\left(x_{0}\right)=f\left(x_{0}\right), p^{\prime}\left(x_{1}\right)=f^{\prime}\left(x_{1}\right), p\left(x_{2}\right)=f\left(x_{2}\right) .
\]

若存在，则求出所有的 \(p(x)\) ，若不存在，则给出一个反例。

\subsection{题目4}
四．（16 分）设 \(A, B\) 分别为 \(m \times n\) 和 \(n \times m\) 矩阵，\(I_{n}\) 为 \(n\) 阶单位矩阵。证明：

\[
r(A)+r(B)=r(A B)+n
\]

的充要条件是存在矩阵 \(X, Y\) ，使得 \(X A-B Y=I_{n}\) ．

\subsection{题目5}
五．（16 分）设 \(\mathbb{C}^{n \times n}\) 是 \(n\) 阶复方阵全体在通常运算下构成的复数域 \(\mathbb{C}\) 上的线性空间，方阵

\[
F=\left(\begin{array}{cccccc}
0 & 0 & 0 & \cdots & 0 & -a_{n} \\
1 & 0 & 0 & \cdots & 0 & -a_{n-1} \\
0 & 1 & 0 & \cdots & 0 & -a_{n-2} \\
\vdots & \vdots & \vdots & & \vdots & \vdots \\
0 & 0 & 0 & \cdots & 0 & -a_{2} \\
0 & 0 & 0 & \cdots & 1 & -a_{1}
\end{array}\right)
\]

（1）假设 \(A=\left(a_{i j}\right)_{n \times n}\) ，且 \(A F=F A\) ，证明：\(A=a_{n 1} F^{n-1}+a_{n-1,1} F^{n-2}+\cdots+a_{21} F+a_{11} E\) 。\\
（2）求 \(\mathbb{C}^{n \times n}\) 的子空间 \(C(F)=\left\{X \in \mathbb{C}^{n \times n} \mid F X=X F\right\}\) 的维数。

\subsection{题目6}
六。（16 分）设 \(\varphi\) 是 \(n\) 维复线性空间上的线性变换，\(f(x)=x^{m}+a_{1} x^{m-1}+\cdots+a_{m}\) 与 \(g(x)\) 是两个复多项式，\(\sigma=f(\varphi), \tau=g(\varphi)\) ．矩阵 \(F\) 是 \(f(x)\) 的友矩阵．已知 \(g(F)\) 是可逆矩阵．证明：

\[
\operatorname{Ker}(\sigma \tau)=\operatorname{Ker} \sigma \oplus \operatorname{Ker} \tau
\]

\subsection{题目7}
七．（16 分）设 \(V\) 是复数域 \(\mathbb{C}\) 上的 \(n(n>0)\) 维线性空间，\(f, g\) 是 \(V\) 上的线性变换，满足 \(f g-g f=f\)。证明：\(f\) 的特征值都是 0 ，且．\(f, g\) 有公共的特征向量．

\subsection{题目8}
八．（18分）设 \(A\) 为 \(n\) 阶实方阵，试找出 \(\lim _{k \rightarrow \infty} A^{k}\) 存在的充分必要条件．

\subsection{题目9}
九．（16 分）设 \(S\) 为实对称方阵，证明：当且仅当 \(S\) 可逆时，存在方阵 \(A\) ，使得 \(S A+A^{\mathrm{T}} S\) 正定．
\end{document}
